
\section{\hard\ \trRegular\trLocalRing{}的进阶事实}

正则局部环几乎具备你想要的每一种好性质，但其中一些需要下苦功。我们现在来讨论几条更为精致的事实，它们或许能帮你夜里睡得安稳。

\textbf{13.8.1.\ 正则局部环的诸\trLocalization{}都是\trRegular\trLocalRing。}

\textbf{13.8.2.\ 事实（[E, Cor.~19.14]、[Mat2, Thm.~19.3]）。}\ ——\ \emph{若 $(A,\mathfrak{m})$ 是一个\trRegular\trLocalRing，则 $A$ 在任一\trPrimeIdeal{}处的\trLocalization{}也是\trRegular\trLocalRing。}

（我们并不需要这条结果，因此也不会去证明它。）这条重大定理在\trCommutative{}代数中曾是一个长期悬而未决的问题，直到 \trSerre{}与 Auslander--Buchsbaum 用同调方法将其解决。

于是，要检验 $\Spec A$ 是否正则（$A$ \trNoetherian），只需在诸\trClosedPoint{}（即诸\trMaximalIdeal{}）处检验就够了。假定事实~13.8.2（并使用\mbox{练习~5.1.E}），你便可以通过在闭点处检验来检验一个\trNoetherian{}概形的正则性。

我们将证明事实~13.8.2 的两个重要特例。第一个你现在就能动手做。

\textbf{13.8.A.（练习）}\quad 设 $X$ 是一个\trNoetherian{}、维数为 $1$ 的概形，并且在其诸\trClosedPoint{}处正则。证明：$X$ 是\trReduced{}的。进而（在不援引事实~13.8.2 的前提下）证明：$X$ 是正则的。

第二个重要特例将在 \S21.6.2 中得到证明：

\textbf{13.8.3.\ \trhTheorem{}。}\ ——\ \emph{若 $X$ 是完全域 $k$ 上的一个\trFiniteType{}概形，并且在其诸\trClosedPoint{}处正则，则 $X$ 正则。}

更一般地，\mbox{练习~21.6.B}将证明：若 $A$ 是完全域上一个有限型代数的\trLocalization，则事实~13.8.2 成立。

\textbf{13.8.B.（练习：推广\mbox{练习~13.3.E}）}\quad 设 $k$ 是特征为 $0$ 的\trAlgClosedField。假定定理~13.8.3，证明：$\mathbb{P}^n$ 中存在每个正次数 $d$ 的正则\trHypersurface。（与\mbox{练习~13.3.E}一样，尽管随意放宽假设。）

\textbf{13.8.4.\ 正则局部环都是\trUFD、整闭的，并且是 \trCM{}的。}

\textbf{13.8.5.\ 事实（Auslander--Buchsbaum 定理）。}\ ——\ \emph{正则局部环都是\trUFD。}

（这是一条艰难的定理，所以我们不证明它，因而也不会去使用它。证明可参看 [E, Thm.~19.19]、[Mat2, Thm.~20.3] 或 [Gr-EGA, IV$_4$.21.11.1]；至于\trVariety{}这一特殊情形，则见 [Mu7, \S III.7]。）从而正则概形都是\trFactorial{}的，故由\mbox{练习~5.4.F}是\trNormal{}的。

\textbf{13.8.6.}\ \emph{\trhRemark{}：\trFactoriality{}弱于\trRegularity。}\quad 「正则蕴涵唯一分解」这条蕴含关系是严格的。下面就是一个表明此事的例子。设 $k$ 是特征不为 $2$ 的\trAlgClosedField。令 $A=k[x_1,\dots,x_n]/(x_1^2+\cdots+x_n^2)$（对照\mbox{练习~5.4.I}）。注意 $\Spec A$ 显然在原点处\trSingular。在\mbox{练习~15.5.R}中，我们将证明当 $n\ge5$ 时 $A$ 是一个\trUFD，故 $\Spec A$ 是\trFactorial{}的。特别地，$A_{(x_1,\dots,x_n)}$ 是一个\trNoetherian\trLocalRing，它是\trUFD，却不是正则的。（更一般地，\trGrothendieck{}对 P.~Samuel 一个猜想的证明蕴涵：一个\trNoetherian\trLocalRing，倘若它是一个\trCompleteIntersection（特别地，是一个\trHypersurface），并且在余维数至多 $3$ 处\trFactorial，则它必是\trFactorial{}的，[SGA2, Exp.~XI, Cor.~3.14]。一份更简短且多少更初等的证明见 [CL]。）

\textbf{13.8.7.}\ \emph{正则局部环是整闭的。}\quad Auslander--Buchsbaum 定理（事实~13.8.5）蕴涵：正则局部环是整闭的（由\mbox{练习~5.4.F}）。我们将在 \S26.3.5 中证明这件事（不诉诸事实~13.8.5）。

\textbf{13.8.8.}\ \emph{正则局部环是 \trCM{}的。}\quad 在 \S26.2.5 中，我们将证明正则局部环都是 \trCM{}的（这个概念将在第 26 章中定义）。
